%% Методические указания
%% Лаборатория математики ФН1
\documentclass[12pt, a4paper]{article}
\usepackage[T2A]{fontenc}
\usepackage[utf8]{inputenc}
\usepackage[russian]{babel}
\usepackage[left=25mm, right=20mm, top=20mm, bottom=20mm]{geometry}
\usepackage{amsmath, amssymb}
\usepackage{enumitem}
\pagestyle{plain}
\setlength{\parindent}{1.25cm}
\begin{document}
\begin{center}
  {\large\bfseries Методические указания}

  \vspace{0.3em}
  {\large Решение систем линейных уравнений методом Гаусса}

  \vspace{0.5em}
  {\small Линейная алгебра · 1 курс\\ Лаборатория математики ФН1}
\end{center}
\vspace{1em}

\section*{Порядок действий}

\begin{enumerate}[nosep]
  \item Выпишите расширенную матрицу системы $(A\,|\,b)$.
  \item Прямой ход: элементарными преобразованиями строк приведите матрицу
        к ступенчатому виду. Меняйте строки местами, чтобы ведущий элемент не был нулём.
  \item Сравните ранги: $\operatorname{rk} A$ и $\operatorname{rk}(A\,|\,b)$.
        \begin{itemize}[nosep]
          \item ранги различны~--- система несовместна, решений нет;
          \item ранги равны числу неизвестных~--- единственное решение;
          \item ранги равны, но меньше числа неизвестных~--- бесконечно много решений.
        \end{itemize}
  \item Обратный ход: выразите базисные неизвестные через свободные.
  \item Для неопределённой системы запишите общее решение и фундаментальную
        систему решений однородной системы.
  \item \textbf{Проверка обязательна:} подставьте найденное решение в исходную систему.
\end{enumerate}

\section*{Типичные ошибки}

\begin{itemize}[nosep]
  \item Складывают столбцы вместо строк. Элементарные преобразования системы~--- только по строкам.
  \item Забывают преобразовывать столбец свободных членов.
  \item Получив ступенчатый вид, объявляют систему определённой, не сравнив ранг с числом неизвестных.
  \item В общем решении оставляют базисные неизвестные справа~--- выражать нужно их через свободные.
  \item Не выполняют проверку. Она занимает минуту и ловит большинство арифметических ошибок.
\end{itemize}

\section*{О параметрах}

Если в системе есть параметр, ведущий элемент может обращаться в нуль.
Разбирайте случаи отдельно: сначала общий (параметр не равен особым значениям),
затем каждое особое значение. В типовом расчёте потеря случая~--- самая частая
причина снижения балла.

\end{document}
